% CC BY 4.0 License
% Copyright (c) 2026 Nanum Learn

\documentclass[a4paper,11pt]{article}

\usepackage[T1]{fontenc}
\usepackage[utf8]{inputenc}
\usepackage[english]{babel}

\usepackage{amsmath}
\usepackage{amssymb}
\usepackage{amsthm}
\usepackage{mathtools}

\usepackage[margin=1in]{geometry}
\usepackage{lastpage}
\usepackage{fancyhdr}
\fancyhead[L]{NL01CHEM01 Session 1 Problem Set Solutions}
\fancyhead[R]{Nanum Learn}
\fancyfoot[C]{}
\pagestyle{fancy}
\usepackage[skip=10pt]{parskip}
\setlength{\headheight}{13.59999pt}
\setlength{\headsep}{8pt}

\usepackage{titlesec}
\titleformat{\section}
    {\normalfont\LARGE\bfseries\color{BrickRed}}
    {\S\! \thesection}
    {0.5em}
    {}
\titleformat{\subsection}
    {\normalfont\Large\bfseries}
    {\thesubsection}
    {0.5em}
    {}

\usepackage{enumitem}
\usepackage[dvipsnames]{xcolor}


\newcommand{\easy}{(\ensuremath{\bullet \circ \circ})}
\newcommand{\medium}{(\ensuremath{\bullet \bullet \circ})}
\newcommand{\hard}{(\ensuremath{\bullet \bullet \bullet})}

\newtheoremstyle{listprobstyle}
    {\parsep}{\parsep}
    {\normalfont}
    {}
    {\bfseries}
    {}
    { }
    {{\thmname{#1}~\thmnumber{#2}.}~
        \thmnote{\textcolor{BrickRed}{\bfseries #3}}\\}
\theoremstyle{listprobstyle}
\newtheorem{problem}{Problem}
\newenvironment{solution}
    {\par\begingroup\color{BrickRed}}
    {\par\endgroup}



\title{\textbf{NL01CHEM01 Session 1 Problem Set Solutions}}
\date{3 November 2026}
\author{William Kang}

\begin{document}

\maketitle
\begin{center}
    \small\textit{Edited and Formatted by Enoch Yu}
\end{center}

\section{Solutions}
\begin{problem}[Element, Compound, or Mixture? \easy]
Write E, C, or M next to each.
\vspace{-0.6\baselineskip}
\begin{enumerate}[label=\alph*),itemsep=0pt,parsep=0pt,topsep=0pt]
\item Copper wire (\textcolor{BrickRed}{E})
\begin{solution}
Copper is one element.
\end{solution}
\vspace{0.5\baselineskip}

\item Table salt (\textcolor{BrickRed}{C})
\begin{solution}
Salt is the compound sodium chloride.
\end{solution}
\vspace{0.5\baselineskip}

\item Well-stirred salt water (\textcolor{BrickRed}{M})
\begin{solution}
Saltwater is a mixture.
\end{solution}
\vspace{0.5\baselineskip}

\item Room air (\textcolor{BrickRed}{M})
\begin{solution}
Air is mostly nitrogen and oxygen mixed.
\end{solution}
\vspace{0.5\baselineskip}

\item Distilled water (\textcolor{BrickRed}{C})
\begin{solution}
Distilled water is the compound $\mathrm{H_2O}$.
\end{solution}
\end{enumerate}
\end{problem}


\begin{problem}[Pure, Uniform, or Non-Uniform? \easy]
Label each as PURE, UNIFORM, or NON-UNIFORM.
\vspace{-0.6\baselineskip}
\begin{enumerate}[label=\alph*),itemsep=0pt,parsep=0pt,topsep=0pt]
\item Spoonful of stirred sugar water:
    \underline{\textcolor{BrickRed}{UNIFORM}}
\begin{solution}
Stirred sugar water looks the same throughout but has sugar + water.
\end{solution}
\vspace{0.5\baselineskip}

\item Handful of trail mix:
    \underline{\textcolor{BrickRed}{NON-UNIFORM}}
\begin{solution}
Trail mix shows visible parts.
\end{solution}
\vspace{0.5\baselineskip}

\item Sealed distilled water:
    \underline{\textcolor{BrickRed}{PURE}}
\begin{solution}
Distilled water has fixed composition.
\end{solution}
\vspace{0.5\baselineskip}

\item Rested Italian dressing:
    \underline{\textcolor{BrickRed}{NON-UNIFORM}}
\begin{solution}
Rested dressing separates.
\end{solution}
\end{enumerate}
\end{problem}

\begin{problem}[A Mystery Block \medium]
A small solid block has mass $54\,\mathrm{g}$ and volume
$20\,\mathrm{cm^3}$. (a) Compute its density. (b) Which of the
following elements is the solid block most likely made of:
aluminum ($2.7\,\mathrm{g/cm^3}$), iron ($7.9\,\mathrm{g/cm^3}$),
or copper ($8.9\,\mathrm{g/cm^3}$)?
\end{problem}
\vspace{-\baselineskip}
\begin{solution}
Notice that $\rho = 54/20 = 2.7\,\mathrm{g/cm^3}$. Therefore,
the block is most likely aluminum. Although color, shape, and size
were not available, the measured density narrows down the choices.
\end{solution}

\begin{problem}[Counting Parts of an Atom \medium]
An atom has atomic number 17 and mass number 35.
\vspace{-0.6\baselineskip}
\begin{enumerate}[label=\alph*),itemsep=0pt,parsep=0pt,topsep=0pt]
\item Which element is it?
\begin{solution}
chlorine ($Z = 17$)
\end{solution}
\vspace{0.5\baselineskip}

\item Number of protons?
\begin{solution}
17 protons
\end{solution}
\vspace{0.5\baselineskip}

\item Number of neutrons?
\begin{solution}
18 neutrons ($A - Z = 35 - 17$)
\end{solution}
\vspace{0.5\baselineskip}

\item Number of electrons if neutral?
\begin{solution}
17 electrons for neutral atom
\end{solution}
\vspace{0.5\baselineskip}
\end{enumerate}
\end{problem}


\begin{problem}[Average Atomic Mass \hard]
An element has two isotopes. Isotope M weighs $10.0\,\mathrm{u}$
and makes up 20.0\% of atoms. Isotope N weighs $11.0\,\mathrm{u}$
and makes up 80.0\%. Find the average atomic mass to one decimal
place.
\end{problem}
\vspace{-\baselineskip}
\begin{solution}
The approximate average atomic mass can be represented as
$(10.0)(0.20) + (11.0)(0.80) = 2.0 + 8.8 = 10.8\,\mathrm{u}$.
This is because when multiplying each isotope mass by its fraction and
adding the products, the total fractions sum must to 1.00.
\end{solution}


\begin{problem}[Glass of Water \hard]
Use $18\,\mathrm{g/mol}$ for water and $6.02 \times 10^{23}$
particles/mol.
\vspace{-0.6\baselineskip}
\begin{enumerate}[label=\alph*),itemsep=0pt,parsep=0pt,topsep=0pt]
\item How many moles in $180\,\mathrm{g}$?
\begin{solution}
$180/18 = 10\,\mathrm{mol}$
\end{solution}
\vspace{0.5\baselineskip}

\item How many molecules total $\rightarrow$ write in scientific
notation.
\begin{solution}
$10 \times 6.02 \times 10^{23} = 6.02 \times 10^{24}$ molecules
\end{solution}
\vspace{0.5\baselineskip}

\item Roughly how many per sip if a glass is six sips?
\begin{solution}
About $1 \times 10^{24}$ per sip
\end{solution}
\vspace{0.5\baselineskip}
\end{enumerate}
\end{problem}
\end{document}


